% !TeX root = ../深度学习与大语言模型技术原理.tex
% 由主文档通过 \input 载入；请编译主文档。
\setcounter{chapter}{18} % 保留原书章号。
\chapter{离线评测、评分器与在线实验}
\label{ch:evaluation}

\section{先定义任务，再定义指标}

评测首先明确输入分布、允许信息、成功条件和预算。选择题、自由生成、代码执行和环境交互的成功定义不同；同一模型在不同提示模板、工具权限和采样次数下也不是同一个被评测系统。报告应固定模型及分词器版本、提示、解码设置、上下文限制、工具和评分器版本。

基准测试（benchmark）提供可复用任务集合，但受训练污染、难度饱和和分布代表性限制。精确匹配需要规范化约定，过强规范化可能把不同答案合并；选择题通过选项似然还是直接生成字母评分，也会产生不同结果。不能把不同协议的分数当成同一实验比较。

错误样本（bad case）分析应形成可操作类别，例如任务理解、知识缺失、证据检索、计算、格式、工具参数、环境执行和评分错误。每次修复应对应类别和可验证假设，并用新验证数据确认；不断查看测试样本并修补会使测试集退化成开发集。

\section{成功率、配对比较与不确定性}

若 $N$ 个独立任务的成功指标为 $Z_i\in\{0,1\}$，经验成功率 $\widehat p=N^{-1}\sum_iZ_i$。在同分布 Bernoulli 假设下，其标准误可估为 $\sqrt{\widehat p(1-\widehat p)/N}$；样本少或接近边界时，应采用更适当的区间方法，不能机械使用正态近似。来自同一模板或同一用户的任务可能相关，应按组重采样或聚类处理。

比较模型 A、B 时，最好在相同任务上计算差值 $D_i=Z_i^A-Z_i^B$，再估计 $\overline D$ 的不确定性。这样利用了题目难度的配对信息。多个随机种子、解码随机性和多次调参带来的选择偏差应明确记录；只报告表现最好的一次运行会高估预期能力。

代码生成常用 pass@$k$。对同一道题独立生成 $n$ 个候选，其中 $c$ 个通过，$n\ge k$，无放回随机取 $k$ 个候选至少一次通过的概率估计为
\begin{equation}
\widehat{\operatorname{pass@}k}=1-\frac{\binom{n-c}{k}}{\binom nk},
\label{eq:passk}
\end{equation}
其中 $n-c<k$ 时分子取零。该估计针对同一采样机制下的候选成功概率\cite{humaneval}；它假设能辨认哪个候选通过，不能代替实际单答案选择系统的准确率。

\section{评分器工程与验证器的可靠性}

评分器（grader）将输出转成分数或标签，验证器（verifier）检验某种条件是否满足。规则评分适合精确结果、格式和单元测试；模型评分适合难以规则化的语义判断；混合评分结合外部执行与语义评估。无论哪种方式，都应把评分器作为独立待验证组件。

答案提取应保留原始输出，记录提取失败与数学错误的区别。数值比较要规定绝对和相对容差，符号等价要说明定义域与异常点；程序测试要固定依赖、输入和资源限制。模型输出若截断或工具环境损坏，应进入单独错误类别，不能无声计成正确或排除出分母。

评分器测试集应覆盖正确、明显错误、边界和诱导样本，估计误接受率和误拒绝率。对训练奖励而言，误接受尤其容易被策略主动利用；对评测而言，两种误差都影响排名。应保留独立人工抽检和真实环境验证，并监测训练后评分器失效模式是否变化。

\section{大模型裁判的偏差}

大模型作为裁判（LLM-as-a-judge）可以做成对偏好或按评分表打分。常见偏差包括答案位置、篇幅、语言风格、自我偏好和提示敏感性\cite{judge}。因此可随机交换候选顺序、隐藏模型身份、明确维度，并在人工标注子集上测量一致性。两次交换评判仍不一致时，应保留分歧，不能任意选择有利结果。

裁判与被测模型若共享训练数据或能力缺陷，其错误会相关；增加同源裁判数量不一定等于增加独立证据。对数学、代码和检索引用等可部分验证的任务，应优先结合确定性检查或实际执行。裁判给出的自然语言理由也可能是事后解释，不能单凭理由流畅就认定判定可靠。

\section{在线实验与安全发布}

A/B 测试随机分配用户或请求到两个方案，比较事先规定的任务指标、延迟和成本。随机化单位要与干扰范围一致：同一用户跨请求共享记忆时，按单请求分配可能互相污染。需要预先说明样本量、观测窗口和停止规则，反复看显著性再临时停止会增加误判。

影子部署（shadow deployment）将真实请求复制给新系统观察结果，但不把新系统响应交付用户；有副作用的工具应被隔离或模拟。金丝雀发布（canary release）让小比例流量实际使用新版本，再根据指标逐步调整。二者目的不同，不能把影子流量中的模拟成功视为真实用户收益。

在线指标可以包含任务成功率、人工接管率、用户反馈、尾部延迟和单任务成本。点击、留存或转化率可能受界面和人群影响，并不天然等于答案正确性。故应把离线能力证据与在线行为结果结合，并保留回滚版本和明确故障界限。

